% Study sheet: Volume I, Chapter 5, Neural Computation.
% Numbers: \val{...} from nn_computation_studysheet_numbers.tex, written by nn_computation_studysheet_numbers.py.
\documentclass{studysheet}
\input{nn_computation_studysheet_numbers}
\sheetvolume{1}
\sheetchapter{5}{Neural Computation}
\sheetsources{Numbers from mlsysim and the chapter's scenarios (FP32; 1 KB = 1,000 bytes); assumed inputs are marked.}

\begin{document}
% ===================================================================== FRONT
\frontside
\begin{bigidea}
A neural network replaces hand-written rules with \hl{arithmetic}: matrix products and simple
activations whose shapes fix the work, the bytes and the memory before any code runs. Training
keeps \hl{state that serving never needs}, and it fails when numbers leave their range, not when logic breaks.
\end{bigidea}

\begin{leftcol}
\sheetsec{The running network: a multilayer perceptron}
{\centering
\begin{tikzpicture}[sheet]
  \node[accbox, minimum width=44pt, minimum height=18pt, align=center] (a) at (24,40) {\val{inputs}\\pixels};
  \node[nbox, minimum width=34pt, minimum height=18pt] (b) at (100,40) {\val{nh1}};
  \node[nbox, minimum width=34pt, minimum height=18pt] (c) at (165,40) {\val{nh2}};
  \node[accbox, minimum width=40pt, minimum height=18pt, align=center] (d) at (228,40) {\val{nout}\\scores};
  \draw[arr] (a) -- (b); \draw[arr] (b) -- (c); \draw[arr] (c) -- (d);
  \node[small, anchor=south] at (66,46) {$\val{inputs}{\times}\val{nh1}$};
  \node[small, anchor=south] at (133,46) {$\val{nh1}{\times}\val{nh2}$};
  \node[small, anchor=south] at (196,46) {$\val{nh2}{\times}\val{nout}$};
  % MACs per layer, to scale
  \pgfmathsetmacro\sc{246}
  \pgfmathsetmacro\xa{2+\val{f1}*\sc} \pgfmathsetmacro\xb{\xa+\val{f2}*\sc} \pgfmathsetmacro\xc{\xb+\val{f3}*\sc}
  \draw[hotbox] (2,0) rectangle (\xa,14);
  \node[lab, text=white] at ({(2+\xa)/2},7) {layer 1: \val{w1} MACs (\val{l1share}\%)};
  \draw[nbox] (\xa,0) rectangle (\xb,14);
  \draw[nbox] (\xb,0) rectangle (\xc,14);
  \node[small, anchor=north east] at (\xc,-3) {layer 2: \val{w2}; layer 3: \val{w3}};
\end{tikzpicture}\par}
\vspace{5pt}
{\fontsize{8.6}{11}\selectfont\color{muted}Each hidden unit: a weighted sum of all its inputs plus a bias, then ReLU. Softmax turns the
\val{nout} scores into probabilities. \val{weights} weights $=$ \val{macs} MACs per digit; with \val{biases} biases,
\val{params} parameters, \val{wkb} KB.\par}

\vspace{8pt}
\fcard{Counting a dense network (per example)}%
  {\Parameters = \sum_{\LayerIndex} (\PreviousLayerWidth + 1)\,\CurrentLayerWidth \qquad
   \Operations = 2\sum_{\LayerIndex} \PreviousLayerWidth\,\CurrentLayerWidth}%
  {Parameters $\Parameters$: a weight per input per unit, plus a bias per unit; $\CurrentLayerWidth$ is layer $\LayerIndex$'s width.
   Bytes $=$ count $\times$ \val{bpv} in FP32.}%
  {}{}

\vspace{12pt}
\sheetsec{Key ideas}
\begin{ideas}
\idea{A layer is a matrix product plus an activation.}{$\PreactivationMatrix = \ActivationMatrix\WeightMatrix + \BiasVector$,
  then $\ActivationFunction$, one row per example. Without $\ActivationFunction$, the layers collapse into one affine map.}
\idea{Shapes fix the cost.}{One multiply-accumulate (MAC, 2 FLOPs) per weight per example. Widening adjacent
  layers costs quadratically; depth, linearly but in sequence.}
\idea{Batching reuses weights.}{Batch \val{batch} shares each weight read: about \val{i32} FLOP/byte, not \val{i1}.
  Tiny networks may still be ruled by per-operation overhead.}
\idea{Learning fails by number range.}{Sigmoid's slope (at most \val{sigslope}) shrinks gradients layer by layer; ReLU,
  a cheap comparator, passes 1. Overflow gives Inf or NaN.}
\idea{A threshold turns scores into decisions.}{Items below it go to people. Pick the lowest
  $\ManualReviewFraction\ManualReviewCost + \MisclassifiedFraction\ErroneousPredictionCost$, where $\MisclassifiedFraction =$ accepted share $\times$ error among accepted.}
\end{ideas}
\end{leftcol}%
\begin{rightcol}
\sheetsec{Formulas you will use}
\fcard{Backward pass of one dense layer}%
  {\frac{\partial \LossFunction}{\partial \WeightMatrix^{(\LayerIndex)}} = {\ActivationMatrix^{(\LayerIndex-1)}}^{T} \frac{\partial \LossFunction}{\partial \PreactivationMatrix^{(\LayerIndex)}}
   \qquad \frac{\partial \LossFunction}{\partial \ActivationMatrix^{(\LayerIndex-1)}} = \frac{\partial \LossFunction}{\partial \PreactivationMatrix^{(\LayerIndex)}} {\WeightMatrix^{(\LayerIndex)}}^{T}}%
  {Two forward-sized products: one reads the saved input, one carries the signal back. Layer 1 skips
   the second, so a step costs about 3 forward passes, less here.}%
  {Then every parameter steps against its gradient: $\ModelParameters \leftarrow \ModelParameters - \LearningRate\,\nabla_{\ModelParameters}\LossFunction$.}{}
\fcard{Training against inference memory}%
  {\begin{aligned}\text{training} &\approx \text{parameters} + \text{gradients} + \text{Adam} + \text{saved}\\
   \text{inference} &\approx \text{parameters} + \text{widest pair of outputs}\end{aligned}}%
  {Gradients match the parameters; Adam keeps two values per parameter. Training saves every layer's output,
   input included: batch $\BatchSize \times \sum \CurrentLayerWidth$ values.}%
  {Inference frees each output once the next layer reads it: $\BatchSize(\PreviousLayerWidth + \CurrentLayerWidth)$ values at the widest pair.}{}
\fcard{Cross-entropy and its gradient}%
  {\LossFunction = -\ln \CorrectClassProbability \qquad
   \frac{\partial \LossFunction}{\partial \NeuronPreactivation_{\ClassIndex}} = \PredictedOutputComponent_{\ClassIndex} - \TrueLabelComponent_{\ClassIndex}}%
  {Gradient: predicted minus target (1 for the true class,
   else 0); a negative one raises that score.}%
  {A uniform guess costs \val{q2uni}. Compute the loss fused from max-shifted scores, so no probability rounds to zero before the log.}{}
\end{rightcol}

% ===================================================================== BACK
\backside
\begin{answers}
\ans{Finish}{Activations \val{bakb} KB; training \val{statekb} $+$ \val{bakb} $=$ \val{btrain} KB. Inference
  \val{wkb} $+$ \val{bikb} $=$ \val{binfer} KB. Ratio \val{bratio}$\times$: saved activations are now
  \val{bactshare}\% of training memory, so the batch, not the parameters, sets it. [Memory card]}
\ans{1}{Weights \val{q1w1} $+$ \val{q1w2} $+$ \val{q1w3} $=$ \val{q1w} $=$ MACs per digit (\val{q1x}$\times$
  the running network); $+$ \val{q1b} biases $=$ \val{q1p} parameters, \val{q1mb} MB. The first layer holds \val{q1share}\%. [Counting card]}
\ans{2}{$-\log \val{q2pbad} = \val{q2lbad}$ against $-\log \val{q2pok} = \val{q2lok}$, about \val{q2lx}$\times$.
  Gradients: true digit \val{q2pbad} $-$ \val{q2target} $=$ \val{q2gtrue} (push up); wrong digit $+$\val{q2gwrong} (push down). [Cross-entropy card]}
\columnbreak
\ans{3}{A: vanishing gradients; the sigmoid slopes alone allow at most \val{sigslope}$^{\val{sigdepth}} \approx \val{sigsurv}$
  by the first layer. Switch to ReLU. B: a range failure (overflow, or the log of a probability rounded to zero); compute
  the loss fused from the scores. [range idea; cross-entropy card]}
\ans{4}{Per \val{q4per} items, A: \val{q4an} reviews $\times$ \$\val{q4cr} $+$ \val{q4aw} errors $\times$ \$\val{q4ce}
  $=$ \$\val{q4acost}. B: \val{q4bn} $\times$ \$\val{q4cr} $+$ \val{q4bw} $\times$ \$\val{q4ce} $=$ \$\val{q4bcost}.
  B is cheaper: its extra reviews cost less than A's extra errors. [threshold idea]}
\end{answers}

\begin{leftcol}
\sheetsec{Worked example}
\begin{practicebox}\boxtitle{How much more memory does training need than serving?}
The running network, batch \val{batch}, FP32, trained with Adam.\par\vspace{5pt}
{\centering
\begin{tikzpicture}[sheet]
  \pgfmathsetmacro\sc{150/\val{trainkbraw}}
  \pgfmathsetmacro\pa{50+\val{wkbraw}*\sc} \pgfmathsetmacro\pb{\pa+\val{gkbraw}*\sc}
  \pgfmathsetmacro\pc{\pb+\val{okbraw}*\sc} \pgfmathsetmacro\pd{\pc+\val{akbraw}*\sc}
  \pgfmathsetmacro\qb{\pa+\val{ikbraw}*\sc}
  \node[note, anchor=west] at (0,31) {training};
  \draw[nbox] (50,24) rectangle (\pa,38); \node[lab] at ({(50+\pa)/2},31) {params};
  \draw[nbox] (\pa,24) rectangle (\pb,38); \node[lab] at ({(\pa+\pb)/2},31) {grads};
  \draw[nbox] (\pb,24) rectangle (\pc,38); \node[lab] at ({(\pb+\pc)/2},31) {Adam};
  \draw[hotbox] (\pc,24) rectangle (\pd,38);
  \node[small, anchor=south] at ({(\pc+\pd)/2},40) {saved};
  \node[lab, anchor=west] at ({\pd+4},31) {\val{trainkb} KB};
  \node[note, anchor=west] at (0,8) {inference};
  \draw[nbox] (50,1) rectangle (\pa,15); \node[lab] at ({(50+\pa)/2},8) {params};
  \draw[accbox] (\pa,1) rectangle (\qb,15);
  \node[small, anchor=south] at ({(\pa+\qb)/2},17) {buffers};
  \node[lab, anchor=west] at ({\qb+4},8) {\val{inferkb} KB};
\end{tikzpicture}\par}
\begin{steps}
\item Parameters: \val{params} $\times$ \val{bpv} B $=$ \val{wkb} KB. Gradients \val{gkb} KB; Adam \val{okb} KB:
  \res{\val{statekb} KB} of parameter state.
\item Saved: \val{batch} $\times$ (\val{inputs}$+$\val{nh1}$+$\val{nh2}$+$\val{nout}) $=$ \val{actvals} values $=$ \val{akb} KB. Training \res{\val{trainkb} KB}.
\item Inference: \val{batch} $\times$ (\val{inputs}$+$\val{nh1}) $=$ \val{infvals} values $=$ \val{ikb} KB, plus parameters:
  \res{\val{inferkb} KB}, so training needs \res{\val{memratio}$\times$}.
\end{steps}
\takeaway{At small batch, the parameter state, not the activations, makes training dearer.}
\end{practicebox}

\vspace{12pt}
\sheetsec{Your turn: finish it}
\begin{practicebox}\boxtitle{Raise the batch to \val{bbatch}}
Same network, \val{bgrow}$\times$ the batch.\par\vspace{4pt}
\textbf{Given:} parameter state is unchanged, \val{statekb} KB. Saved values: \val{bbatch} $\times$ \val{actsum}
$=$ \val{bactvals}. Inference buffers: \val{bbatch} $\times$ \val{pairw} $=$ \val{binfvals} values $=$ \val{bikb} KB.\par
\textbf{You:} find training and inference memory, their ratio, and which term now dominates training.
\end{practicebox}
\end{leftcol}%
\begin{rightcol}
\sheetsec{Common mistakes}
\mistake{``A data-center accelerator will train this network far faster than my CPU.''}%
  {Modeled: its math and memory traffic per epoch run \val{epspeed}$\times$ faster (\val{epworkcpu} to \val{epworkacc} ms),
   but each of \val{epops} small operations per step pays a fixed launch cost: \val{eptotacc} ms against
   the CPU's \val{eptotcpu} ms, \val{epovh}\% of it overhead.}{}
\mistake{``The backward pass runs the network in reverse, so it costs one forward pass.''}%
  {Each layer needs two products of forward size, weight gradient and input gradient. This network's step
   costs \val{trainratio}$\times$ the forward work, not 3$\times$, only because layer 1 needs no gradient for the raw pixels.}{}
\mistake{``The 1989 ZIP-code reader had \val{ziperr}\% error, so machines sorted 99\% of the mail.''}%
  {The \val{ziperr}\% held only for the digits it accepted, after sending \val{ziprej}\% to people.}{}

\vspace{2pt}
\sheetsec{Check yourself}
\begin{qs}
\item \assumed{} a \val{inputs}$\to$\val{q1h1}$\to$\val{q1h2}$\to$\val{nout} network in FP32. Count its parameters,
  MACs per digit and megabytes. Which layer dominates?
\item \assumed{} for a seven, the network gives probability \val{q2pbad} to ``seven'' and \val{q2pwrong} to ``three''. Find the loss,
  compare it with \val{q2pok} on ``seven'', and give the gradient at both scores.
\item A \val{sigdepth}-layer sigmoid network stalls. In run A the loss is flat and gradients shrink toward the first layers;
  in run B the loss becomes NaN. Which failure does each suggest, and what would you change first?
\item \assumed{} review costs \$\val{q4cr}, an accepted error \$\val{q4ce}. Threshold A reviews \val{q4arev}\%
  and errs on \val{q4aerr}\% of accepted items; B reviews \val{q4brev}\% and errs on \val{q4berr}\%. Which is cheaper?
\end{qs}
\end{rightcol}
\end{document}
